% Part 3 of Day 8: Deployment.
%
% Topics of the syllabus: input resolution and model scale; export formats;
% int8 calibration for detectors; latency of the full pipeline: capture,
% pre-process, infer, post-process. Exercise: predict the frame rate change
% when the resolution goes from 640 to 320.
%
% The measured numbers of this part are results of an experiment of this
% course: YOLO11n of Ultralytics, its exported files, and 500 images of the
% validation set of COCO, on the work computer (an x86 processor with AVX2,
% 4 threads). The times say nothing about a board. Each frame with such a
% number says so in its source line. The times for a Raspberry Pi are
% numbers that the author of the companion book published. The frame names
% the source.

\theorypart{Deployment}%
  {predict the frame rate when the resolution goes from 640 to 320}

% ---------------------------------------------------------------- decisions
\begin{frame}{From a trained model to a frame rate}
  \begingroup\centering
  \begin{tikzpicture}[font=\scriptsize, >={Stealth[length=4pt]},
      b/.style={draw=coolgray, rounded corners=3pt, fill=boxgray,
                minimum width=1.75cm, minimum height=0.95cm, align=center}]
    \node[b] (a) at (0,0)    {Trained\\model};
    \node[b, draw=darkcyan]    (b) at (2.2,0) {Image size\\and scale};
    \node[b, draw=darkorange]  (c) at (4.4,0) {Export\\format};
    \node[b, draw=cardinalred] (d) at (6.6,0) {Precision};
    \node[b, draw=treegreen]   (e) at (8.8,0) {Pipeline on\\the board};
    \foreach \a/\b in {a/b, b/c, c/d, d/e} { \draw[->] (\a) -- (\b); }
    \node[coolgray] at (2.2,-0.8) {320 or 640};
    \node[coolgray] at (4.4,-0.8) {NCNN, LiteRT};
    \node[coolgray] at (6.6,-0.8) {\texttt{float32}, \texttt{int8}};
    \node[coolgray] at (8.8,-0.8) {camera to result};
  \end{tikzpicture}\par\endgroup
  \vskip 0.4em
  \small
  \begin{itemize}
    \item Part 2 gave the accuracy of a model. A product also needs a frame
          rate: the detections of each second.
    \item Four decisions set the frame rate. Each decision can cost
          accuracy.
    \item The syllabus of the lab asks for a table: two image sizes and two
          precisions, with the frame rate and the accuracy of each row.
    \item This part gives the method for each decision, with measured
          values of YOLO11n.
  \end{itemize}
  \keypoint{Report the frame rate and the accuracy together, for the
    file that runs on the board.}
\end{frame}

% --------------------------------------------------------------- resolution
\begin{frame}{The cost grows with the pixels}
  \begin{columns}[c]
    \begin{column}{0.29\textwidth}
      \centering
      \begin{tikzpicture}[font=\scriptsize, x=0.0044cm, y=0.0044cm]
        \fill[darkcyan!12] (0,0) rectangle (640,640);
        \draw[darkcyan] (0,0) rectangle (640,640);
        \fill[darkcyan!26] (0,0) rectangle (480,480);
        \draw[darkcyan] (0,0) rectangle (480,480);
        \fill[darkcyan!42] (0,0) rectangle (320,320);
        \draw[darkcyan] (0,0) rectangle (320,320);
        \fill[darkcyan!60] (0,0) rectangle (224,224);
        \draw[darkcyan] (0,0) rectangle (224,224);
        \node[anchor=north east] at (640,640) {640};
        \node[anchor=north east] at (480,480) {480};
        \node[anchor=north east] at (320,320) {320};
        \node[anchor=north east, white] at (224,224) {224};
      \end{tikzpicture}
    \end{column}
    \begin{column}{0.68\textwidth}
      \small
      \renewcommand{\arraystretch}{1.2}
      \setlength{\tabcolsep}{5pt}
      \begin{tabular}{@{}rrrr@{}}
        \toprule
        \textbf{Image side} & \textbf{Pixels} & \textbf{Candidates}
        & \textbf{GFLOP} \\
        \midrule
        640 & 100 percent & 8400 & 6.5 \\
        480 & 56 percent  & 4725 & 3.7 \\
        320 & 25 percent  & 2100 & 1.6 \\
        224 & 12 percent  & 1029 & 0.8 \\
        \bottomrule
      \end{tabular}
    \end{column}
  \end{columns}
  \vskip 0.4em
  \small
  \begin{itemize}
    \item Each convolution does its work for each pixel of its map. One
          half of the image side gives one quarter of the operations.
    \item The parameters do not change: the file has the same size.
    \item Measured for the network in ONNX Runtime: 49.6 ms at 640 pixels
          and 13.6 ms at 320 pixels. The factor is 3.6.
  \end{itemize}
  \keypoint{The image size is the strongest control of the latency. It
    needs no new model.}
  \source{An experiment of this course on an x86 processor, 4 threads.}
\end{frame}

\begin{frame}{Measured: the resolution against the accuracy}
  \begin{columns}[c]
    \begin{column}{0.50\textwidth}
      \centering
      \begin{tikzpicture}[font=\scriptsize, x=1.15cm, y=6.0cm]
        \foreach \k/\a/\b/\l in {0/0.574/0.417/640, 1/0.535/0.391/480,
                                 2/0.451/0.319/320, 3/0.353/0.244/224} {
          \fill[darkcyan!35] (\k,0) rectangle (\k+0.38,\a);
          \fill[darkcyan!85] (\k+0.42,0) rectangle (\k+0.80,\b);
          \node[above, font=\tiny] at (\k+0.19,\a) {\a};
          \node[above, font=\tiny] at (\k+0.61,\b) {\b};
          \node[below] at (\k+0.4,0) {\l};
        }
        \draw[coolgray] (-0.15,0) -- (4,0);
        \node[coolgray] at (1.9,-0.1) {image side in pixels};
      \end{tikzpicture}
    \end{column}
    \begin{column}{0.48\textwidth}
      \small
      \begin{itemize}
        \item Light bars: mAP50. Dark bars: mAP50-95.
        \item The same weights, four image sizes. The model was trained
              at 640 pixels.
        \item From 640 to 480 pixels: 44 percent fewer operations, and
              mAP50-95 falls by 0.026.
        \item From 640 to 320 pixels: 75 percent fewer operations, and
              mAP50-95 falls by 0.099.
      \end{itemize}
    \end{column}
  \end{columns}
  \vskip 0.3em
  \keypoint{A smaller image is not free. Measure the mAP at each size that
    you want to use.}
  \source{YOLO11n on 500 images of COCO, an experiment of this course.}
\end{frame}

\begin{frame}{Small objects pay first}
  \begingroup\centering
  \begin{tikzpicture}[font=\scriptsize, x=2.35cm, y=3.3cm]
    \foreach \k/\a/\b/\c/\d/\l in {
        0/0.23/0.16/0.07/0.01/{Small objects},
        1/0.54/0.46/0.31/0.17/{Medium objects},
        2/0.79/0.77/0.72/0.62/{Large objects}} {
      \fill[darkcyan!90] (\k,0)      rectangle (\k+0.19,\a);
      \fill[darkcyan!65] (\k+0.21,0) rectangle (\k+0.40,\b);
      \fill[darkcyan!42] (\k+0.42,0) rectangle (\k+0.61,\c);
      \fill[darkcyan!22] (\k+0.63,0) rectangle (\k+0.82,\d);
      \node[above, font=\tiny] at (\k+0.095,\a) {\a};
      \node[above, font=\tiny] at (\k+0.305,\b) {\b};
      \node[above, font=\tiny] at (\k+0.515,\c) {\c};
      \node[above, font=\tiny] at (\k+0.725,\d) {\d};
      \node[below] at (\k+0.41,0) {\l};
    }
    \draw[coolgray] (-0.1,0) -- (3.0,0);
    \node[anchor=west, align=left] at (3.05,0.52)
      {Recall at\\640, 480, 320,\\and 224 pixels,\\from dark to light};
  \end{tikzpicture}\par\endgroup
  \vskip 0.2em
  \small
  \begin{itemize}
    \item Small: a label with an area below $32 \times 32$ pixels in the
          photo. Large: above $96 \times 96$ pixels.
    \item At 320 pixels, the model finds 7 percent of the small objects.
          It still finds 72 percent of the large objects.
    \item An object of 32 pixels covers 4 cells on a side of the finest
          map. At one half of the image size, it covers 2 cells.
  \end{itemize}
  \keypoint{Ask how many pixels the smallest object of your product has.
    This number selects the image size.}
  \source{YOLO11n on 500 images of COCO: score threshold 0.25, IoU
    threshold 0.5. An experiment of this course.}
\end{frame}

% -------------------------------------------------------------------- scale
\begin{frame}{The model scale: n, s, m, l, x}
  \small
  \renewcommand{\arraystretch}{1.2}
  \begin{tabular}{@{}lrrrr@{}}
    \toprule
    \textbf{Model} & \textbf{Parameters} & \textbf{GFLOP}
    & \textbf{Published} & \textbf{500 images} \\
    \midrule
    YOLO11n & 2.6 million  & 6.5   & 0.395 & 0.417 \\
    YOLO11s & 9.4 million  & 21.6  & 0.470 & 0.503 \\
    YOLO11m & 20.1 million & 68.1  & 0.515 & 0.551 \\
    YOLO11l & 25.3 million & 87.2  & 0.534 & not measured \\
    YOLO11x & 56.9 million & 195.3 & 0.547 & not measured \\
    \bottomrule
  \end{tabular}
  \vskip 0.4em
  \begin{itemize}
    \item The two last columns give mAP50-95 at 640 pixels. A scale changes
          the number of channels and the number of layers.
    \item From n to s: 3.3 times the operations for 0.086 more mAP50-95.
          From s to m: 3.2 times the operations for 0.047 more.
    \item The gain of each step gets smaller. The cost of each step does
          not.
    \item The Raspberry Pi 5 has no accelerator. The lab uses the scale n.
  \end{itemize}
  \keypoint{Start with the smallest scale. Go up only when the accuracy is
    not sufficient and the frame rate has a margin.}
  \source{Published: the documentation of Ultralytics, 5000 images of
    COCO. Measured: an experiment of this course.}
\end{frame}

\begin{frame}{A larger model, or a larger image?}
  \begin{columns}[c]
    \begin{column}{0.56\textwidth}
      \centering
      \begin{tikzpicture}[font=\scriptsize, x=2.25cm, y=11cm,
          >={Stealth[length=4pt]}]
        \draw[->, coolgray] (-0.2,0.2) -- (2.05,0.2);
        \draw[->, coolgray] (-0.2,0.2) -- (-0.2,0.6);
        \node[coolgray, anchor=west] at (-0.18,0.59) {mAP50-95};
        \node[coolgray] at (0.9,0.155) {GFLOP};
        \foreach \x/\l in {0/1, 1/10, 2/100} {
          \draw[coolgray] (\x,0.2) -- (\x,0.193) node[below] {\l};
        }
        \foreach \y in {0.3,0.4,0.5} {
          \draw[coolgray] (-0.2,\y) -- (-0.22,\y) node[left] {\y};
          \draw[coolgray!30, dotted] (-0.2,\y) -- (2.0,\y);
        }
        \draw[darkcyan, thick] (-0.100,0.244) -- (0.211,0.319) --
          (0.564,0.391) -- (0.816,0.417);
        \foreach \x/\y in {-0.100/0.244, 0.211/0.319, 0.564/0.391,
                           0.816/0.417} {
          \fill[darkcyan] (\x,\y) circle (1.8pt);
        }
        \draw[cardinalred, thick] (0.730,0.407) -- (1.334,0.503);
        \fill[cardinalred] (0.730,0.407) circle (1.8pt);
        \fill[cardinalred] (1.334,0.503) circle (1.8pt);
        \fill[darkorange] (1.833,0.551) circle (1.8pt);
        \node[darkcyan!60!black, anchor=north west] at (-0.08,0.243)
          {n, 224};
        \node[darkcyan!60!black, anchor=south east] at (0.82,0.42)
          {n, 640};
        \node[cardinalred, anchor=north west] at (0.74,0.405) {s, 320};
        \node[cardinalred, anchor=south east] at (1.33,0.505) {s, 640};
        \node[darkorange, anchor=north] at (1.833,0.545) {m, 640};
      \end{tikzpicture}
    \end{column}
    \begin{column}{0.42\textwidth}
      \small
      \begin{itemize}
        \item The horizontal axis has a logarithmic scale.
        \item YOLO11n at 640 pixels: 6.5 GFLOP and 0.417.
        \item YOLO11s at 320 pixels: 5.4 GFLOP and 0.407.
        \item The two choices cost almost the same, and they give almost
              the same mAP.
        \item The small image of YOLO11s loses the small objects. Its file
              is 3.4 times as large.
      \end{itemize}
    \end{column}
  \end{columns}
  \vskip 0.3em
  \keypoint{The mAP does not decide here. The size of your objects and the
    memory of the device decide.}
  \source{500 images of COCO, an experiment of this course.}
\end{frame}

% ---------------------------------------------------------------- size rule
\begin{frame}{The image size rule: train, export, run}
  \begingroup\centering
  \begin{tikzpicture}[font=\footnotesize, >={Stealth[length=5pt]},
      b/.style={draw=coolgray, rounded corners=3pt, fill=boxgray,
                minimum width=2.4cm, minimum height=0.95cm, align=center}]
    \node[b, draw=darkcyan] (a) at (0,0)   {Train\\\texttt{imgsz=320}};
    \node[b, draw=darkcyan] (b) at (3.6,0) {Export\\\texttt{imgsz=320}};
    \node[b, draw=darkcyan] (c) at (7.2,0) {Run\\\texttt{imgsz=320}};
    \draw[->] (a) -- (b); \draw[->] (b) -- (c);
  \end{tikzpicture}\par\endgroup
  \vskip 0.5em
  \small
  \begin{itemize}
    \item The PyTorch file accepts each image size. An exported file has a
          static input shape (Day 7): the size is a part of the file.
    \item The export uses 640 pixels when you give no size.
    \item Published case: a model that was trained at 320 pixels and
          exported with no size gave 113 boxes and 164 wheels for an image
          with one box. The program gave no error and no warning.
    \item A model that was trained at 320 pixels saw objects of a certain
          size in pixels. At a different size, the objects look different
          to it.
  \end{itemize}
  \keypoint{Use one image size for the training, the export, and the
    inference. Check the shape in the first output line.}
  \source{\emph{Edge AI Engineering}, ``Computer Vision Applications with
    YOLO''.}
\end{frame}

% ---------------------------------------------------------------- letterbox
\begin{frame}{Pre-processing: the letterbox}
  \begingroup\centering
  \begin{tikzpicture}[font=\scriptsize, >={Stealth[length=4pt]},
      x=0.0019cm, y=0.0019cm]
    \fill[darkcyan!25] (0,0) rectangle (1920,1080);
    \draw[darkcyan] (0,0) rectangle (1920,1080);
    \node at (960,540) {camera frame: $1920 \times 1080$};
    \draw[->, thick, coolgray] (2020,540) -- (2420,540);
    \begin{scope}[xshift=4.9cm, x=0.0039cm, y=0.0039cm, yshift=-0.22cm]
      \fill[coolgray!35] (0,0) rectangle (640,640);
      \fill[darkcyan!25] (0,140) rectangle (640,500);
      \draw[darkcyan] (0,0) rectangle (640,640);
      \node at (320,320) {$640 \times 360$};
      \node at (320,570) {border};
      \node at (320,70) {border};
    \end{scope}
  \end{tikzpicture}\par\endgroup
  \vskip 0.2em
  \small
  \begin{itemize}
    \item The program makes the frame smaller with one factor for the two
          sides, and grey borders fill the square. After the inference, it
          moves each box back to the frame.
    \item For a frame of 16:9, the borders are 44 percent of the square.
          The network does this work for nothing.
    \item An export with \code{imgsz=384,640} gives 5040 candidates, not
          8400. Measured in ONNX Runtime: 29.9 ms, not 49.6 ms.
  \end{itemize}
  \keypoint{The pre-processing of the device must be the pre-processing of
    the training. A different resize changes the accuracy.}
  \source{An experiment of this course on an x86 processor.}
\end{frame}

% ------------------------------------------------------------------ formats
\begin{frame}{Export formats}
  \small
  \renewcommand{\arraystretch}{1.2}
  \setlength{\tabcolsep}{4pt}
  \begin{tabular}{@{}L{0.20\textwidth}L{0.29\textwidth}rrL{0.20\textwidth}@{}}
    \toprule
    \textbf{Format} & \textbf{Files} & \textbf{Size} & \textbf{Export}
    & \textbf{Runtime} \\
    \midrule
    PyTorch & \code{yolo11n.pt}                & 5.35 MB  & none
            & PyTorch \\
    ONNX    & \code{yolo11n.onnx}              & 10.24 MB & 5 s
            & ONNX Runtime \\
    NCNN    & \code{.param} and \code{.bin}    & 10.10 MB & 8 s
            & NCNN \\
    LiteRT  & \code{yolo11n.tflite}            & 10.27 MB & 19 s
            & LiteRT \\
    LiteRT, \code{int8} & \code{yolo11n\_int8.tflite} & 2.94 MB & 136 s
            & LiteRT \\
    \bottomrule
  \end{tabular}
  \vskip 0.4em
  \begin{itemize}
    \item The PyTorch file stores the weights in \code{float16}. The three
          exports store \code{float32}: the files are two times as large.
    \item The \code{int8} export needs calibration images. It ran the model
          on 200 images.
    \item The three runtimes are the runtimes of Day 7. The export fuses
          each batch normalization into its convolution.
  \end{itemize}
  \source{YOLO11n at 640 pixels, an experiment of this course. 1 MB is
    1024 KB.}
\end{frame}

\begin{frame}[fragile]{Export and run with one tool}
\begin{lstlisting}[style=shell]
# On the laptop: export the trained model for the two runtimes
yolo export model=best.pt format=ncnn imgsz=320
yolo export model=best.pt format=litert imgsz=320 quantize=8 data=data.yaml

# On the Raspberry Pi: run, then measure
yolo predict model=best_ncnn_model imgsz=320 source=test.jpg
yolo val model=best_int8.tflite imgsz=320 data=data.yaml
\end{lstlisting}
  \small
  \begin{itemize}
    \item \code{best.pt} is the file of your training. The export writes
          the folder \code{best\_ncnn\_model} and the file
          \code{best\_int8.tflite}.
    \item \code{quantize=8} makes a model with \code{int8} weights and
          \code{int8} activations. \code{data} names the dataset: the
          export takes the calibration images from its validation set.
    \item The export to LiteRT runs on a laptop with Linux or macOS. It
          does not run on the Raspberry Pi.
    \item \code{predict} and \code{val} accept an exported file. The
          package selects the runtime from the name.
  \end{itemize}
  \source{The Ultralytics package, version 8.4.171. The export commands
    ran on the work computer with \texttt{yolo11n.pt} at 640 pixels.}
\end{frame}

\begin{frame}{Measured: one model in four files}
  \small
  \renewcommand{\arraystretch}{1.25}
  \begin{tabular}{@{}L{0.30\textwidth}rrr@{}}
    \toprule
    \textbf{File and runtime} & \textbf{mAP50} & \textbf{mAP50-95}
    & \textbf{Network, 4 threads} \\
    \midrule
    PyTorch                    & 0.579 & 0.424 & 69.1 ms \\
    ONNX in ONNX Runtime       & 0.579 & 0.424 & 49.6 ms \\
    NCNN                       & 0.579 & 0.424 & 60.9 ms \\
    LiteRT, \code{float32}     & 0.579 & 0.424 & 34.8 ms \\
    \bottomrule
  \end{tabular}
  \vskip 0.5em
  \begin{itemize}
    \item The four files give the same mAP. An export in \code{float32}
          does not change the accuracy.
    \item The times are from an x86 processor. NCNN is made for Arm: on a
          Raspberry Pi, the order can change.
    \item This test puts each image into a square of 640 pixels. With the
          rectangle of Part 2, the same model has 0.417. The
          pre-processing is a part of the result.
  \end{itemize}
  \keypoint{After each export, run the validation again with the exported
    file. Equal values are the proof.}
  \source{YOLO11n on 500 images of COCO, an experiment of this course on
    an x86 processor. The times say nothing about a board.}
\end{frame}

% --------------------------------------------------------------------- int8
\begin{frame}{int8 for a detector: boxes and scores}
  \begingroup\centering
  \begin{tikzpicture}[font=\scriptsize, >={Stealth[length=4pt]}]
    \draw[coolgray, thick] (0,0.9) -- (8.5,0.9);
    \foreach \x in {0,8.5} { \draw[coolgray, thick] (\x,0.8) -- (\x,1.0); }
    \node[left] at (0,0.9) {0};
    \node[right] at (8.5,0.9) {640};
    \node[darkcyan!60!black] at (4.25,1.2)
      {a box value in pixels: one step of 8 bits is 2.5 pixels};
    \fill[cardinalred] (0,0.82) rectangle (0.035,0.98);
    \draw[->, cardinalred] (0.6,0.35) -- (0.06,0.8);
    \node[cardinalred, anchor=west] at (0.65,0.3)
      {all scores are between 0 and 1: they fall into one step};
  \end{tikzpicture}\par\endgroup
  \vskip 0.3em
  \small
  \begin{itemize}
    \item The output tensor of YOLO holds box values from 0 to 640 and
          scores from 0 to 1. A tensor in \code{int8} has one scale
          (Day 4).
    \item One scale for 640 gives steps of 2.5. Each score becomes 0: the
          model finds nothing.
    \item The export of the package divides the box values by the image
          size. All values are then between 0 and 1.
    \item In the \code{int8} file of the experiment, the scale of the
          output is 0.0039. One step of a box value is 2.5 pixels at 640
          pixels.
    \item The input and the output of the file stay in \code{float32}. The
          file converts at its two ends.
  \end{itemize}
  \keypoint{Look at the range of each output before you quantize a
    detector. Two ranges in one tensor need a change of the model.}
  \source{The export code of the Ultralytics package, and the exported
    file read with LiteRT. An experiment of this course.}
\end{frame}

\begin{frame}{Measured: int8 and the calibration}
  \footnotesize
  \renewcommand{\arraystretch}{1.25}
  \setlength{\tabcolsep}{4pt}
  \begin{tabular}{@{}L{0.39\textwidth}rrrr@{}}
    \toprule
    \textbf{LiteRT file} & \textbf{Size} & \textbf{mAP50}
    & \textbf{mAP50-95} & \textbf{Network} \\
    \midrule
    \code{float32}                           & 10.27 MB & 0.579 & 0.424
        & 34.8 ms \\
    \code{int8} weights only                 & 2.90 MB  & 0.580 & 0.419
        & 53.7 ms \\
    \code{int8}, calibration: 200 photos     & 2.94 MB  & 0.517 & 0.331
        & 55.5 ms \\
    \code{int8}, calibration: noise          & 2.94 MB  & 0.493 & 0.299
        & not measured \\
    \code{int8} weights, \code{int16} activations & 3.00 MB & 0.573
        & 0.415 & 2329 ms \\
    \bottomrule
  \end{tabular}
  \vskip 0.3em
  \small
  \begin{itemize}
    \item \code{int8} weights alone cost 0.005 of mAP50-95. Full
          \code{int8} costs 0.062 of mAP50 and 0.093 of mAP50-95: the
          boxes lose more than the classes.
    \item Calibration with noise costs 0.032 more of mAP50-95. Use photos
          of your application, as on Day 4.
    \item With \code{int16} activations, the loss is 0.009 of mAP50-95.
          But LiteRT has no fast kernels for them on this processor: the
          network needs 67 times the time of \code{float32}.
    \item On this x86 processor, \code{int8} is slower than
          \code{float32}. For a classifier on the Raspberry Pi 5, a
          published test gave \code{int8} a factor of 2.9 (Day 7).
  \end{itemize}
  \source{An experiment of this course on an x86 processor.}
\end{frame}

% ---------------------------------------------------------------- published
\begin{frame}{Published: times on Raspberry Pi boards}
  \small
  \renewcommand{\arraystretch}{1.25}
  \setlength{\tabcolsep}{4pt}
  \begin{tabular}{@{}L{0.34\textwidth}L{0.22\textwidth}rr@{}}
    \toprule
    \textbf{Model} & \textbf{Board} & \textbf{PyTorch} & \textbf{NCNN} \\
    \midrule
    YOLO11n for box and wheel, 320 pixels & Raspberry Pi 5
        & about 400 ms & about 80 ms \\
    YOLO11n for COCO, 640 pixels          & Raspberry Pi Zero 2 W
        & 4672 ms & 2748 ms \\
    \bottomrule
  \end{tabular}
  \vskip 0.5em
  \begin{itemize}
    \item The times are for the network alone.
    \item On the Raspberry Pi 5, the export to NCNN gives a factor of 5.
          It costs no accuracy.
    \item The SSD MobileNet V1 file of Part 1 needs about 100 ms on a
          Raspberry Pi 5 in LiteRT.
    \item The sources give no time for a YOLO model in LiteRT with
          \code{int8} on the Raspberry Pi 5. You measure it in the lab.
  \end{itemize}
  \keypoint{First export to a runtime for the board. Then measure.}
  \source{\emph{Edge AI Engineering}, chapters on detection and YOLO.}
\end{frame}

% ----------------------------------------------------------------- pipeline
\begin{frame}{The pipeline of one frame}
  \begingroup\centering
  \begin{tikzpicture}[font=\scriptsize, >={Stealth[length=4pt]},
      b/.style={draw=coolgray, rounded corners=3pt, fill=boxgray,
                minimum width=1.75cm, minimum height=0.95cm, align=center}]
    \node[b, draw=darkcyan]    (a) at (0,0)   {\textbf{Capture}};
    \node[b, draw=darkorange]  (b) at (2.2,0) {\textbf{Pre-process}};
    \node[b, draw=cardinalred] (c) at (4.4,0) {\textbf{Infer}};
    \node[b, draw=treegreen]   (d) at (6.6,0) {\textbf{Post-process}};
    \node[b] (e) at (8.8,0) {Use the\\result};
    \foreach \a/\b in {a/b, b/c, c/d, d/e} { \draw[->] (\a) -- (\b); }
  \end{tikzpicture}\par\endgroup
  \vskip 0.4em
  \small
  \renewcommand{\arraystretch}{1.2}
  \begin{tabular}{@{}L{0.17\textwidth}L{0.43\textwidth}L{0.32\textwidth}@{}}
    \toprule
    \textbf{Step} & \textbf{Work} & \textbf{Who does it} \\
    \midrule
    Capture      & Wait for the next frame of the camera. Copy it to the
                   program.
                 & The camera and its driver \\
    Pre-process  & Letterbox, colour order, \code{float32}, values from 0
                   to 1
                 & OpenCV and numpy on the CPU \\
    Infer        & The network & The runtime, with its threads \\
    Post-process & Score threshold, suppression, boxes back to the frame
                 & Python code on one core \\
    Use the result & Draw, count, send a message (Week 3)
                 & Your program \\
    \bottomrule
  \end{tabular}
  \vskip 0.3em
  \keypoint{The frame rate is one divided by the sum of all steps, not by
    the time of the network.}
\end{frame}

\begin{frame}{The time of each step: two devices}
  \begingroup\centering
  \begin{tikzpicture}[font=\scriptsize, x=0.064cm, y=0.62cm]
    \foreach \k/\a/\b/\c/\l in {
        0/3.8/93.4/2.8/{x86 computer, 640 pixels},
        1/18.8/41.1/40.1/{Raspberry Pi Zero 2 W, 320 pixels}} {
      \fill[darkorange!60]  (0,-\k) rectangle (\a,-\k+0.7);
      \fill[cardinalred!60] (\a,-\k) rectangle (\a+\b,-\k+0.7);
      \fill[treegreen!60]   (\a+\b,-\k) rectangle (100,-\k+0.7);
      \node[anchor=east] at (-1,-\k+0.35) {\l};
    }
    \node at (50,0.35) {network: 93 percent};
    \node at (9.4,-0.65) {19};
    \node at (39.3,-0.65) {network: 41 percent};
    \node at (80,-0.65) {post-process: 40};
    \node[darkorange!70!black, anchor=west] at (0,1.1) {pre-process};
    \node[treegreen!60!black, anchor=east] at (100,1.1) {post-process};
  \end{tikzpicture}\par\endgroup
  \vskip 0.3em
  \footnotesize
  \renewcommand{\arraystretch}{1.15}
  \setlength{\tabcolsep}{5pt}
  \begin{tabular}{@{}L{0.33\textwidth}rrrr@{}}
    \toprule
    & \textbf{Pre-process} & \textbf{Network} & \textbf{Post-process}
    & \textbf{Frame rate} \\
    \midrule
    x86 computer, measured           & 2.0 ms   & 49.3 ms  & 1.5 ms
        & 19.0 \\
    Raspberry Pi Zero 2 W, published & 237.6 ms & 519.3 ms & 506.7 ms
        & 0.8 \\
    \bottomrule
  \end{tabular}
  \vskip 0.3em
  \small
  \begin{itemize}
    \item The frame rate is in frames each second. On the small board, the
          network alone would give 1.9.
    \item The pre-processing and the post-processing are Python code. They
          get no gain from a faster runtime or from \code{int8}.
    \item Published for the same board: 2545 ms for the network in the
          first run, and 519 ms in the second run. Do not count the first
          run.
  \end{itemize}
  \source{x86: an experiment of this course, YOLO11n in ONNX Runtime.
    Raspberry Pi: \emph{Edge AI Engineering}, chapter on YOLO.}
\end{frame}

\begin{frame}{What limits the frame rate}
  \small
  \renewcommand{\arraystretch}{1.2}
  \begin{tabular}{@{}L{0.24\textwidth}L{0.68\textwidth}@{}}
    \toprule
    \textbf{Limit} & \textbf{What happens} \\
    \midrule
    The camera     & A camera with 30 frames each second gives one frame
                     each 33 ms. A faster pipeline waits. \\
    The sum of the steps & A pipeline of 100 ms gives 10 results each
                     second. \\
    Old frames     & The driver keeps frames in a queue. A slow program
                     reads frames that are old. \\
    The cores      & Four threads for the network leave no core for the
                     capture and for Python (Day 7). \\
    The heat       & At full load, the processor can reduce its clock
                     after some minutes (Day 9). \\
    \bottomrule
  \end{tabular}
  \vskip 0.3em
  \begin{itemize}
    \item Read the camera in a separate thread, and keep only the newest
          frame.
    \item The frame rate says how many results you get. The latency says
          how old a result is.
  \end{itemize}
  \keypoint{Measure the frame rate over one minute of the real pipeline.}
\end{frame}

% ------------------------------------------------------------------ actions
\begin{frame}{How to reach a target frame rate}
  \small
  \renewcommand{\arraystretch}{1.2}
  \begin{tabular}{@{}rL{0.34\textwidth}L{0.52\textwidth}@{}}
    \toprule
    & \textbf{Action} & \textbf{Cost} \\
    \midrule
    1 & Export to a runtime for the board
      & None. The mAP stays the same. \\
    2 & Use a rectangle for a frame of 16:9
      & None. The model does not see the borders. \\
    3 & Use a smaller image
      & The small objects. Measure the mAP again. \\
    4 & Use \code{int8}
      & Some mAP, mostly at the strict thresholds. It needs calibration
        photos. \\
    5 & Use a smaller scale
      & Accuracy, and the time for a new training \\
    6 & Run the network on each second frame
      & A moving object is found later. \\
    \bottomrule
  \end{tabular}
  \vskip 0.4em
  \begin{itemize}
    \item After each action, measure the frame rate of the pipeline on the
          board and the mAP of the file. Stop when the target is reached.
  \end{itemize}
  \keypoint{Take the actions in this order: the first actions cost
    nothing. Write each step in the Decision Log.}
\end{frame}

\begin{frame}{In the lab}
  \small
  \renewcommand{\arraystretch}{1.25}
  \begin{tabular}{@{}lL{0.24\textwidth}L{0.56\textwidth}@{}}
    \toprule
    \textbf{Part} & \textbf{Topic} & \textbf{Work} \\
    \midrule
    A & Pre-trained models
      & Run the SSD model and a YOLO model on test images on the
        Raspberry Pi. \\
    B & Custom model
      & Train a YOLO model on a small dataset with two classes, on the
        laptop or on Colab. \\
    C & Export and deploy
      & Export to NCNN and to LiteRT in \code{int8}. Run the live
        detection with the camera. \\
    D & Measure
      & The frame rate for two image sizes and two precisions. The
        accuracy loss of the \code{int8} model. \\
    \bottomrule
  \end{tabular}
  \vskip 0.5em
  \begin{itemize}
    \item The check: the model detects the two classes in the live image,
          and the table has four measured rows.
  \end{itemize}
  \keypoint{Write your prediction for each row before Part D. Nobody
    measured these four rows on a Raspberry Pi 5 for this course.}
\end{frame}

% -------------------------------------------------------------------- exercise
\exerciseframe{the frame rate from 640 to 320}{%
  \small
  A Raspberry Pi 5 runs a detector on the frames of its camera. The table
  gives example values for an image size of 640 pixels.
  \vskip 0.2em
  {\footnotesize
  \renewcommand{\arraystretch}{1.1}
  \begin{tabular}{@{}L{0.50\textwidth}r@{}}
    \toprule
    Capture and pre-processing & 15 ms \\
    Network                    & 300 ms \\
    Post-processing            & 10 ms \\
    \bottomrule
  \end{tabular}}
  \vskip 0.3em
  \begin{enumerate}
    \item Calculate the frame rate at 640 pixels.
    \item You export the model for 320 pixels. Predict the time of each
          step and the new frame rate.
    \item By which factor does the frame rate change? Why is the factor
          not 4?
    \item An object has $40 \times 40$ pixels in the image of 640 pixels.
          Which size does it have at 320 pixels? What do you expect for
          its recall?
    \item The target is 15 frames each second. Name the next action.
  \end{enumerate}
  Work with your neighbour for 8 minutes.
}

\answerframe{the frame rate from 640 to 320}{%
  \small
  \begin{enumerate}
    \item $15 + 300 + 10 = 325$ ms. The frame rate is $1000 / 325 = 3.1$
          frames each second.
    \item The network has one quarter of the pixels: $300 / 4 = 75$ ms.
          The capture and the post-processing change little: 15 ms and
          10 ms. The sum is 100 ms: 10 frames each second.
    \item The factor is $325 / 100 = 3.25$. The 25 ms outside the network
          do not change. This is the law of Amdahl of Day 7.
    \item $20 \times 20$ pixels. The object changes from ``medium'' to
          ``small''. In the experiment, the recall of small objects is
          far below the recall of medium objects: expect a large loss.
    \item 15 frames each second need 67 ms. The network must go from 75 ms
          to about 42 ms. Try \code{int8}, and measure the mAP. A
          rectangle for a frame of 16:9 costs no accuracy: try it first.
  \end{enumerate}
  \vskip 0.2em
  Measured on the x86 computer: the network went from 49.3 ms to 13.6 ms
  (a factor of 3.6), and the pipeline from 52.8 ms to 15.7 ms (a factor of
  3.4).
}
